\title{DeepSeekMath: Pushing the Limits of Mathematical Reasoning in Open Language Models}

\begin{abstract}
The inherently intricate and structured aspects of mathematical reasoning present considerable obstacles for language models. In this study, we present DeepSeekMath~7B, an extension of DeepSeek-Coder-Base-v1.5 7B, further trained on 120B math-specific tokens gathered from Common Crawl alongside data containing natural language and code. DeepSeekMath~7B attains a notable 51.7\% score on the challenging MATH benchmark at the competition level without utilizing auxiliary toolkits or ensemble voting approaches, thereby nearing the performance of Gemini-Ultra and GPT-4. When using a self-consistency method over 64 generated samples, DeepSeekMath~7B reaches 60.9\% accuracy on MATH. The strengthened mathematical reasoning skills of DeepSeekMath~are largely ascribed to two primary innovations: First, our sophisticated data curation pipeline effectively leverages extensive open-access web resources. Second, we propose Group Relative Policy Optimization (GRPO), an adaptation of Proximal Policy Optimization (PPO), which not only improves the model’s mathematical reasoning but also reduces PPO’s memory footprint.
\end{abstract}

\section{Introduction}

Large language models (LLMs) have transformed mathematical reasoning tasks within artificial intelligence, driving meaningful improvements in benchmarks for quantitative \citep{MATH} and geometric reasoning \citep{trinh2024solving}. Additionally, such models have demonstrated considerable efficacy in aiding humans to tackle intricate mathematical challenges \citep{tao}. Nonetheless, top-performing systems like GPT-4 \citep{gpt4} and Gemini-Ultra \citep{gemini} remain inaccessible to the public, while available open-source models exhibit significant performance gaps by comparison.

This paper presents DeepSeekMath, a domain-focused language model that demonstrably surpasses the mathematical proficiency of current open-source systems, while narrowing the gap with GPT-4 in standardized academic assessments. Central to this progress is the development of the DeepSeekMath Corpus, a substantial and meticulously curated pre-training dataset that encompasses 120B math-specific tokens. The construction of this dataset involves extracting relevant content from the Common Crawl (CC) using a classifier built on fastText \citep{joulin2016fasttext}. In the first phase, the classifier is trained by using samples from OpenWebMath \citep{openwebmath} as positive data and supplementing with a heterogeneous assortment of general web pages for negative samples. The trained classifier is subsequently deployed to retrieve more positive samples from the CC, after which these are rigorously filtered and annotated by human experts. The classifier itself is further refined with this improved dataset, enhancing its extraction accuracy. Empirical analysis demonstrates the robustness of this comprehensive corpus: our foundational model, DeepSeekMath-Base 7B, achieves 64.2\% on GSM8K \citep{gsm8k} and 36.2\% on the advanced-level MATH benchmark \citep{MATH}, thereby exceeding the performance of Minerva 540B \citep{minerva}. Owing to the corpus’s multilingual composition, we also observe improved outcomes in Chinese mathematical evaluation tasks \citep{wei2023cmath,agieval}. Our methodologies in handling and preparing mathematical data serve as a baseline for the wider research community, and there remains considerable potential for continued enhancement moving forward.

DeepSeekMath-Base is constructed by initializing with DeepSeek-Coder-Base-v1.5 7B \citep{deepseek-coder}, as we find that leveraging a model initially trained for coding provides a more advantageous foundation than using a standard large language model. Additionally, we observe that mathematical training not only strengthens the model’s proficiency in math but also yields gains on general benchmarks such as MMLU \citep{mmlu} and BBH \citep{bbh}, thereby improving the model’s overall reasoning capacity beyond mathematics.

Following the pre-training phase, we perform mathematical instruction tuning on DeepSeekMath-Base using data that incorporate chain-of-thought reasoning \citep{cot}, program-of-thought processes \citep{pot,pal}, and tool-assisted reasoning \citep{tora}. The fine-tuned model, DeepSeekMath-Instruct 7B, surpasses all other 7B-scale models and demonstrates performance on par with open-source instruction-tuned models at the 70B scale.

In addition, we present Group Relative Policy Optimization (GRPO), which is a modified reinforcement learning approach inspired by Proximal Policy Optimization (PPO) \citep{schulman2017proximal}. GRPO dispenses with the critic network, instead utilizing group-based score estimations as a baseline, which leads to considerable savings in training resources. Applying GRPO solely with a selected portion of English instruction tuning data yields marked improvements over the already robust DeepSeekMath-Instruct model across both in-domain benchmarks (e.g., GSM8K: 82.9\% $\rightarrow$ 88.2\%, MATH: 46.8\% $\rightarrow$ 51.7\%) and out-of-domain problems (such as CMATH: 84.6\% $\rightarrow$ 88.8\%) during the reinforcement learning process. We also establish a cohesive framework to interpret a variety of related methods, including Rejection Sampling Fine-Tuning (RFT) \citep{yuan2023scaling}, Direct Preference Optimization (DPO) \citep{dpo}, PPO, and GRPO. This unified perspective reveals that these strategies can be regarded as either direct or streamlined forms of reinforcement learning. Through extensive empirical analysis—including comparisons of online versus offline learning, evaluation of outcome versus process-based supervision, and studies contrasting single-turn with iterative RL—we rigorously explore the fundamental aspects of this framework. Finally, we elucidate the mechanisms by which our RL strategies enhance instruction-tuned models and outline prospective avenues for further advancement of RL methods within this integrated paradigm.

\subsection{Contributions}
This work presents advancements in scalable mathematical pre-training and delves into the examination and evaluation of reinforcement learning methodologies.

\textbf{Math Pre-Training at Scale}

\begin{itemize}[topsep=0pt]
    \item We demonstrate that publicly available Common Crawl data possesses substantial utility for mathematical tasks. Through the development of a carefully engineered data selection pipeline, we curate the DeepSeekMath Corpus, comprising 120B tokens of web content curated for mathematical relevance. This dataset notably surpasses previous efforts, being nearly sevenfold larger than the mathematical web page corpus utilized by Minerva \citep{minerva} and approximately nine times larger than the recently released OpenWebMath \citep{openwebmath}.

    \item The DeepSeekMath-Base 7B model, which we pre-trained on this dataset, attains a level of performance on par with Minerva 540B \citep{minerva}. This result suggests that parameter count alone does not determine mathematical reasoning performance. High-quality data and a focused pre-training regimen enable comparatively smaller models to exhibit robust reasoning capabilities.

    \item We report empirical findings from our mathematical training investigations. In particular, pre-training on code prior to math-specific training enhances models’ problem-solving skills, regardless of whether tool usage is involved. This provides evidence supporting a positive relationship between code pre-training and reasoning skills, offering a partial response to the enduring question: \textit{does code training improve reasoning abilities?} At least for mathematical reasoning, our results indicate that it does.

    \item Training using arXiv papers—a practice prevalent in mathematical model development—does not yield significant improvements across the mathematical benchmarks evaluated in this study.
\end{itemize}

\textbf{Exploration and Analysis of Reinforcement Learning}

\begin{itemize}[topsep=0pt]
    \item This work presents Group Relative Policy Optimization (GRPO), a reinforcement learning algorithm designed for both efficiency and effectiveness. Distinct from traditional approaches, GRPO eliminates the need for a critic model by inferring the baseline directly from aggregated group scores. This methodological change greatly decreases the computational demands in comparison to Proximal Policy Optimization (PPO).
    \item We show that GRPO considerably improves the results of our instruction-tuned system, DeepSeekMath-Instruct, leveraging only data from instruction tuning. Moreover, we find that performance gains extend to out-of-domain tasks throughout the reinforcement learning phase.
    \item We propose a comprehensive framework that unifies various techniques, including RFT, DPO, PPO, and GRPO. Our extensive empirical analysis explores a range of aspects such as online versus offline learning, the distinction between outcome and process-based supervision, and comparisons of single-step and iterative reinforcement learning, offering insights into the core components underlying this unified view.
    \item Within this integrative framework, we investigate factors contributing to the efficacy of reinforcement learning and identify promising directions for further advancing reinforcement learning methods applied to LLMs.
\end{itemize}

\subsection{Summary of Evaluations and Metrics}

\begin{itemize}[topsep=0pt]
    \item \textbf{English and Chinese Mathematical Reasoning}: Our evaluation strategy encompasses a broad spectrum of benchmarks in both English and Chinese, ranging from elementary through university-level mathematics. The English benchmarks used are GSM8K \citep{gsm8k}, MATH \citep{MATH}, SAT \citep{llemma}, OCW Courses \citep{minerva}, and MMLU-STEM \citep{mmlu}. For Chinese, the selected benchmarks include MGSM-zh \citep{mgsm}, CMATH \citep{wei2023cmath}, Gaokao-MathCloze \citep{agieval}, and Gaokao-MathQA \citep{agieval}. Our assessment targets not only the capability of producing detailed textual solutions independently—without recourse to external tools—but also evaluates proficiency in solving tasks through Python programming.

    On the English sets, DeepSeekMath-Base performs on par with the proprietary Minerva 540B \citep{minerva}, and it outperforms all other open-source base models such as Mistral 7B \citep{mistral} and Llemma-34B \citep{llemma}, regardless of their pre-training regimens or dataset language composition—often by substantial margins. Particularly in the Chinese benchmarks, DeepSeekMath-Base exhibits clear superiority. This advantage likely arises because, contrary to previous approaches \citep{minerva,llemma}, our methodology does not limit mathematical pre-training data to English, instead incorporating substantial high-quality datasets from other languages. The application of instruction tuning and reinforcement learning further boosts model performance: DeepSeekMath-Instruct and DeepSeekMath-RL each achieve exceptional outcomes, notably surpassing the 50\% accuracy threshold on the rigorous, competition-level MATH dataset—representing a milestone among open-source systems.
\end{itemize}

\item \textbf{Formal Mathematics}: We assess DeepSeekMath-Base on the informal-to-formal theorem proving task as described in \citep{dsp_proof}, utilizing the miniF2F benchmark \citep{minif2f} and selecting Isabelle \citep{isabelle} as the designated proof assistant. The results indicate that DeepSeekMath-Base achieves robust autoformalization performance in a few-shot setting.
\item \textbf{Natural Language Understanding, Reasoning, and Code}: In order to comprehensively evaluate the general abilities of the models in understanding language, reasoning, and programming, DeepSeekMath-Base is benchmarked on the Massive Multitask Language Understanding (MMLU) benchmark \citep{mmlu}—comprised of 57 multiple-choice tests across a wide range of disciplines—as well as BIG-Bench Hard (BBH) \citep{bbh}, featuring 23 tasks that typically require complex multi-step reasoning. Code-specific performance is also measured using HumanEval \citep{codex} and MBPP \citep{mbpp}, both recognized standard benchmarks for code evaluation. Math pre-training is shown to improve results in both language comprehension and reasoning tasks.
\end{itemize}

\section{Math Pre-Training}

\subsection{Data Collection and Decontamination} This section details the procedure for constructing the DeepSeekMath Corpus utilizing Common Crawl data. As illustrated in Figure \ref{fig:data_collect}, our approach employs an iterative pipeline, which begins with a carefully curated seed dataset (for example, a small-scale yet high-quality collection of mathematical content) to methodically expand into a large-scale mathematical dataset from Common Crawl. The framework outlined here is adaptable to other specialized areas, such as programming languages.

Initially, OpenWebMath \citep{openwebmath}, a resource of high-caliber mathematical text from the web, is used as the foundational seed. Leveraging this, we train a fastText model \citep{joulin2016fasttext} in order to retrieve additional web pages with similar mathematical content to OpenWebMath. To do this, we randomly select 500,000 examples from the seed as positive samples and 500,000 unrelated Common Crawl pages as negative examples. The open-source fastText implementation\footnote{\url{https://fasttext.cc}} is used for model training, setting the vector size at 256, a learning rate of 0.1, a maximum word n-gram length of 3, a minimum word occurrence threshold of 3, and running for 3 epochs. To efficiently manage the massive size of Common Crawl, both URL-level deduplication and near-deduplication strategies are implemented, narrowing down the dataset to 40B HTML web pages. Utilizing the trained fastText model, we retrieve mathematically relevant pages from the deduplicated set, and further refine quality by ranking these pages according to their fastText-assigned scores, retaining only those at the upper end of the spectrum. The retained dataset size is validated via pre-training experiments on corpora of the top 40B, 80B, 120B, and 160B tokens. For the initial iteration, the top 40B tokens are selected to proceed.

Following the initial round of data collection, a significant number of mathematical web pages are still omitted, largely due to the insufficiently diverse set of positive examples used for training the fastText model. To enhance the seed corpus and consequently improve model performance, we supplement it with additional mathematical web domains. The Common Crawl dataset is partitioned into non-overlapping domains, defined by grouping all web pages under the same base URL. For each domain, we determine the fraction of pages that were retrieved in the initial round. If more than 10\% of a domain’s web pages are included, that domain is considered math-related (e.g., \url{mathoverflow.net}). Next, we carry out manual annotation to identify URL patterns within these domains that contain mathematical material (for example, \url{mathoverflow.net/questions}). Uncollected web pages corresponding to these URLs are incorporated into the seed corpus. By augmenting the corpus in this way, we acquire a broader set of positive examples, facilitating the training of a more effective fastText model that can identify additional mathematical content in subsequent rounds. After four cycles of data collection, we accumulate 35.5M mathematical web pages, encompassing 120B tokens. During the fourth cycle, we observe that nearly 98\% of the collected data overlaps with the third iteration, leading us to discontinue further data gathering.

To minimize overlap with evaluation datasets (benchmark contamination), we adopt the strategy outlined in \cite{deepseek-coder}, excluding web pages that contain either questions or answers from English mathematical benchmarks such as GSM8K \citep{gsm8k} and MATH \citep{MATH}, as well as Chinese benchmarks like CMATH \citep{wei2023cmath} and AGIEval \citep{agieval}. The filtering procedure involves removing any text segment in our training corpus where a 10-gram sequence matches verbatim with any substring found in these benchmarks. For benchmark snippets that are between 3 and 9 grams, we perform exact matching to eliminate contaminated entries.

\subsection{Evaluating the DeepSeekMath~Corpus Quality}

We conduct pre-training trials to compare the DeepSeekMath~Corpus to recently introduced mathematics-focused corpora:

\begin{itemize}[topsep=0pt]
    \item \textbf{MathPile} \citep{mathpile}: a composite corpus comprising 8.9B tokens, collated from a variety of sources including textbooks, Wikipedia, ProofWiki, CommonCrawl, StackExchange, and arXiv, with arXiv contributing the majority share (over 85\%);
    \item \textbf{OpenWebMath} \citep{openwebmath}: a dataset derived from CommonCrawl through mathematical content filtering, summing to 13.6B tokens;
    \item \textbf{Proof-Pile-2} \citep{llemma}: a collection that combines OpenWebMath, AlgebraicStack (containing 10.3B tokens of mathematical code), and arXiv publications (28.0B tokens). For experiments with Proof-Pile-2, we adhere to the arXiv:Web:Code data ratio of 2:4:1 as prescribed in \cite{llemma}.
\end{itemize}

\subsubsection{Training Setting}
\label{sec:quality-policy}
We conduct math-specific fine-tuning on a foundational language model containing 1.3B parameters, which is architecturally aligned with the DeepSeek LLMs \citep{deepseek-llm}, and is referred to as DeepSeek-LLM 1.3B. Separate models are individually fine-tuned on each mathematical dataset, utilizing 150B training tokens for each case. The entirety of our experimental procedures utilizes the streamlined and efficient HAI-LLM framework \citep{haillm} for training. Adhering to the training methodologies outlined for DeepSeek LLMs, the optimization process uses AdamW \citep{adamW}, setting $\beta_1=0.9$, $\beta_2=0.95$, and $\mathrm{weight\_decay}=0.1$. The learning rate follows a multi-stage decay: beginning with 2,000 warmup steps to reach the maximum, subsequently falling to 31.6\% at 80\% of the training, and further to 10.0\% by 90\% of the training duration. The peak learning rate is assigned as $5.3 \times 10^{-4}$. For each batch, 4M tokens are processed with a context length fixed at 4K.

\subsubsection{Evaluation Results}

\textbf{The DeepSeekMath~Corpus is characterized by superior quality, substantial multilingual representation, and unparalleled scale among existing datasets.}

\begin{itemize}[topsep=0pt]
    \item \textbf{High-quality}:
    Downstream evaluations are carried out across 8 distinct mathematical benchmarks, employing few-shot chain-of-thought prompting \cite{cot}. Table \ref{tab:corpora_comparison} highlights that the model trained with the DeepSeekMath~Corpus consistently achieves top-tier performance. Moreover, Figure \ref{fig:corpora_comparisons} demonstrates that this model outperforms the version trained on Proof-Pile-2 after processing 50B tokens (i.e., one complete epoch of Proof-Pile-2), suggesting that DeepSeekMath~Corpus provides superior data quality on average.
    \item \textbf{Multilingual}:
    The DeepSeekMath~Corpus integrates resources in a variety of languages, with English and Chinese constituting the major share. Empirical results in Table \ref{tab:corpora_comparison} illustrate that training with DeepSeekMath~Corpus boosts mathematical reasoning accuracy in both English and Chinese. By comparison, previous mathematical corpora—predominantly English-focused—deliver limited benefit, or may even adversely affect, performance in Chinese mathematical contexts.
    \item \textbf{Large-scale}:
    DeepSeekMath~Corpus surpasses all current mathematical datasets in terms of total size by a significant margin. Figure \ref{fig:corpora_comparisons} indicates that DeepSeek-LLM 1.3B, when fine-tuned with DeepSeekMath~Corpus, demonstrates accelerated learning progress and continues to improve throughout training. In contrast, baseline corpora are much more limited in volume, leading to repeated cycling during training, and model improvements quickly stagnate.
\end{itemize}

\subsection{Training and Evaluating DeepSeekMath-Base 7B}

This section details the development of DeepSeekMath-Base 7B, a foundational model exhibiting robust mathematical reasoning skills. The model builds upon DeepSeek-Coder-Base-v1.5 7B \citep{deepseek-coder} as its initialization and is further pre-trained over 500B tokens. The dataset composition is distributed as follows: 56\% from the DeepSeekMath Corpus, 4\% from AlgebraicStack, 10\% sourced from arXiv, 20\% consisting of code extracted from Github, and the final 10\% drawn from Common Crawl’s natural language content in both English and Chinese. For training, the general procedures outlined in Section \ref{sec:quality-policy} are largely adhered to, with the primary differences being an increased peak learning rate set at 4.2e-4 and a larger batch size of 10M tokens.

To evaluate the mathematical proficiency of DeepSeekMath-Base 7B, we design a broad suite of tests aimed at three areas: generating fully self-contained mathematical solutions without auxiliary tools, solving mathematics tasks using computational aids, and engaging in formal theorem proving. Additionally, we analyze the broader capabilities of the base model, assessing its competence in natural language processing, deductive reasoning, and programming.

\paragraph{Mathematical Problem Solving with Step-by-Step Reasoning}

DeepSeekMath-Base is assessed on its stepwise reasoning in mathematical problem solving, utilizing few-shot chain-of-thought prompting \citep{cot} across eight diverse benchmarks in both English and Chinese. These benchmarks span a range of quantitative reasoning tasks (including GSM8K \citep{gsm8k}, MATH \citep{MATH}, and CMATH \citep{wei2023cmath}) as well as multiple-choice assessments (such as MMLU-STEM \citep{mmlu} and Gaokao-MathQA \citep{agieval}), covering subject matter from elementary mathematics up to university-level complexity.

As indicated in Table \ref{tab:base_math_cot}, DeepSeekMath-Base 7B consistently achieves top results across all eight benchmark tasks among open-source base models, including the commonly adopted Mistral 7B \citep{mistral} and the more recent Llemma 34B \citep{llemma}, which received additional mathematical training on the Proof-Pile-2 dataset \citep{llemma}. Remarkably, DeepSeekMath-Base exceeds the best open-source competitors by more than 10\% absolute on the MATH dataset—a benchmark reflecting competition-level difficulty—and even surpasses Minerva 540B \citep{minerva}, a proprietary model that is 77 times larger, based on PaLM \citep{palm} and fine-tuned with mathematical material.

\paragraph{Mathematical Problem Solving with Tool Use} 

To assess the ability to perform program-aided mathematical reasoning, we employ few-shot program-of-thought prompting \citep{pot,pal} on the GSM8K and MATH datasets. Each model is prompted to construct a Python program to solve the given question, allowing use of libraries such as \textit{math} and \textit{sympy} for complex computations. The final program output serves as the model's answer. Table \ref{tab:base_math_pot_proof} demonstrates that DeepSeekMath-Base 7B surpasses the previously leading Llemma 34B model in this domain.

\paragraph{Formal Mathematics}

Automating the formalization of mathematical proofs enhances both accuracy and efficiency, a research area gaining considerable traction. DeepSeekMath-Base 7B is assessed on informal-to-formal proof generation as described in \citep{dsp_proof}: given an informal theorem statement, its formal representation, and an informal proof, the task is to synthesize a formal proof. Evaluations utilize miniF2F \citep{minif2f}, a formal Olympiad-level mathematics benchmark. The model generates Isabelle-formatted proofs using few-shot prompting. Consistent with \cite{dsp_proof}, we have models produce proof sketches and subsequently use the Sledgehammer automated prover \citep{sledgehammer} to fill in missing elements. According to Table \ref{tab:base_math_pot_proof}, DeepSeekMath-Base 7B displays strong capabilities for proof autoformalization.

\paragraph{Natural Language Understanding, Reasoning, and Code}

For assessments in natural language understanding, reasoning, and programming, we consider MMLU \citep{mmlu}, BBH \citep{bbh}, HumanEval \citep{codex}, and MBPP \citep{mbpp}, respectively. Table \ref{tab:understanding_reasoning_code} illustrates that DeepSeekMath-Base 7B achieves notable improvements in MMLU and BBH compared with its predecessor DeepSeek-Coder-Base-v1.5 \citep{deepseek-coder}, underscoring the beneficial effects of mathematical training on both language comprehension and reasoning abilities. Furthermore, through the continual inclusion of code tokens during training, DeepSeekMath-Base 7B sustains performance levels on both HumanEval and MBPP, matching DeepSeek-Coder-Base-v1.5. Across all three tasks—reasoning and programming—DeepSeekMath-Base 7B substantially outperforms the generalist Mistral 7B \citep{mistral}.

\section{Supervised Fine-Tuning}

\subsection{SFT Data Curation}

We assembled an instruction-tuning dataset for mathematics, featuring both English and Chinese problems that represent a broad spectrum of mathematical disciplines and difficulty levels. Each problem is paired with a solution, formatted as either chain-of-thought (CoT) \citep{cot}, program-of-thought (PoT) \citep{pot,pal}, or using tool-integrated reasoning approaches \citep{tora}. In total, the dataset comprises 776K training samples.

\begin{itemize}[topsep=0pt]
    \item \textbf{English mathematical datasets}: For English-language resources, we provide tool-integrated solutions for the GSM8K and MATH datasets, utilize a selected subset of MathInstruct \citep{MathInstruct}, and incorporate the Lila-OOD \citep{lila} training set, wherein questions are answered using either CoT or PoT reasoning. This English-language portion spans several mathematical domains, such as algebra, calculus, number theory, probability, and geometry.
    \item \textbf{Chinese mathematical datasets}: In the case of Chinese materials, we collect K-12 mathematics problems from 76 different subfields—including topics like linear equations—with annotations available in both CoT and tool-integrated solution formats.
\end{itemize}

\subsection{Training and Evaluating DeepSeekMath-Instruct 7B}

This section details the supervised tuning process for DeepSeekMath-Instruct 7B, a model developed by fine-tuning DeepSeekMath-Base with mathematically-oriented instructions. We concatenate randomly sampled training items until the token limit of 4K per input is reached. The model is optimized for 500 steps using a batch size of 256 and a fixed learning rate of $5 \times 10^{-5}$.

We measure the model’s proficiency in mathematical tasks, considering both tool-augmented and tool-free scenarios, across four quantitative reasoning benchmarks in both English and Chinese. For context, we compare its performance with contemporary state-of-the-art models:

\begin{itemize}[topsep=0pt]
    \item \textbf{Closed-source models} considered in our comparisons include: (1) the GPT model series, with GPT-4 \citep{gpt4} and GPT-4 Code Interpreter~\footnote{\url{https://openai.com/blog/chatgpt-plugins\#\#code-interpreter}} representing peak performance; (2) Gemini Ultra and Pro \citep{gemini}; (3) Inflection-2 \citep{inflection-2}; (4) Grok-1~\footnote{\url{https://x.ai/model-card}}; as well as several cutting-edge offerings from Chinese organizations such as (5) Baichuan-3~\footnote{\url{https://www.baichuan-ai.com}} and (6) the latest GLM-4~\footnote{\url{https://open.bigmodel.cn/dev/api\#glm-4}} from the GLM family \citep{glm}.
\end{itemize}

The majority of these models are designed for broad application domains, having been subjected to various alignment processes. 
    \item \textbf{Open-source models} span both general-purpose and math-augmented systems. The general models encompass (1) DeepSeek-LLM-Chat 67B \citep{deepseek-llm}, (2) Qwen 72B \citep{qwen}, (3) SeaLLM-v2 7B \citep{seallm}, and (4) ChatGLM3 6B \citep{chatglm3}. Mathematics-optimized models include: (5) InternLM2-Math 20B~\footnote{\url{https://github.com/InternLM/InternLM-Math}}, which extends InternLM2 through targeted math-oriented training and subsequent instruction fine-tuning, (6) Math-Shepherd-Mistral 7B, which incorporates PPO training \citep{schulman2017proximal} into Mistral 7B \citep{mistral} utilizing a process-supervised reward signal, (7) WizardMath series \citep{wizardmath} that upgrades Mistral 7B and Llama-2 70B \citep{llama2} with enhanced mathematical capabilities via evolve-instruct—a variant of instruction tuning employing machine-generated prompts—and PPO, primarily using problems from GSM8K and MATH, (8) MetaMath 70B \citep{metamath}, consisting of Llama-2 70B further refined on expanded GSM8K and MATH datasets, (9) ToRA 34B \cite{tora}, where CodeLlama 34B is fine-tuned for math reasoning augmented with tool use, and (10) MAmmoTH 70B \citep{MathInstruct}, which results from MathInstruct-driven instruction tuning on Llama-2 70B.
\end{itemize}

According to Table \ref{tab:sft_rl_math}, when models are tested in scenarios that restrict external tool usage, DeepSeekMath-Instruct 7B achieves notable proficiency in performing stepwise reasoning. Remarkably, on the competition-grade MATH benchmark, this model exceeds the performance of all open-source systems and most commercial proprietary models (including Inflection-2 and Gemini Pro), offering an absolute margin greater than 9\%. This outperformance holds true even compared to substantially larger models such as Qwen 72B, as well as models that have been especially trained with math-centric reinforcement learning approaches (for instance, WizardMath-v1.1 7B). DeepSeekMath-Instruct matches the results of Chinese proprietary models such as GLM-4 and Baichuan-3 on MATH, although it does not reach the accuracy levels of GPT-4 and Gemini Ultra.

When problem-solving tasks permit the combination of natural language reasoning and programmatic tool integration, DeepSeekMath-Instruct 7B attains nearly 60\% accuracy on the MATH dataset, outstripping all publicly available open-source alternatives. For other evaluation sets, the performance of our model aligns closely with DeepSeek-LLM-Chat 67B, the previous state-of-the-art despite being an order of magnitude larger.

\section{Reinforcement Learning}

\subsection{Group Relative Policy Optimization} Reinforcement learning (RL) techniques have demonstrated their utility in further enhancing the mathematical reasoning skills of large language models after the initial Supervised Fine-Tuning (SFT) stage \citep{wang2023math,wizardmath}. Here, we present our proposed reinforcement learning algorithm, Group Relative Policy Optimization (GRPO), which is both efficient and effective.

\subsubsection{From PPO to GRPO} Proximal Policy Optimization (PPO) \citep{schulman2017proximal} is an actor-critic approach to reinforcement learning, widely adopted for RL-based fine-tuning of large language models \citep{ouyang2022training}. Specifically, PPO seeks to optimize model parameters by maximizing the following surrogate objective:

\begin{equation}
\footnotesize
    \mathcal{J}_{PPO}(\theta) = \mathbb{E}{[q \sim P(Q), o \sim \pi_{\theta_{old}}(O|q)]} \frac{1}{|o|} \sum_{t=1}^{|o|} \min \left[ \frac{\pi_\theta(o_{t} | q, o_{<t})}{\pi_{\theta_{old}}(o_{t} | q, o_{<t})} A_{t}, \text{clip} \left( \frac{\pi_\theta(o_{t} | q, o_{<t})}{\pi_{\theta_{old

Here, $\pi_{\theta}$ and $\pi_{\theta_{old}}$ denote the current and previous policy networks, while $q$ and $o$ represent questions and responses sampled from the question set and from the earlier policy $\pi_{\theta_{old}}$, respectively. The parameter $\varepsilon$ is a clipping hyper-parameter introduced in PPO to help stabilize the learning process. $A_t$ refers to the advantage, which is computed using Generalized Advantage Estimation (GAE) \citep{gae} based on future rewards $\{r_{\ge t}\}$ and an estimated value function $V_{\psi}$. As a consequence, PPO requires simultaneous training of a value function along with the policy; moreover, to control excessive reward model optimization, the standard methodology introduces a per-token KL divergence penalty from a reference model in the reward at every token, as in \citep{ouyang2022training}, specifically,

\begin{equation}
    r_{t} = r_\varphi(q, o_{\le t}) - \beta \log\frac{\pi_{\theta}(o_{t}|q, o_{<t})}{\pi_{ref}(o_{t}|q, o_{<t})},
\label{eq:PPO-reward}
\end{equation}

where $r_\varphi$ is the reward function, $\pi_{ref}$ stands for the reference model (commonly the initial SFT policy), and $\beta$ is the scaling factor for the KL penalty.

Since the value function in PPO is often implemented as a model similar in scale to the policy, this setup can introduce notable computational and memory demands. Furthermore, the value function serves as a baseline during advantage estimation to reduce variance, yet for LLMs, only the reward of the final token is typically provided by the reward model, making token-wise value function training more challenging. To overcome these drawbacks, and as depicted in Figure \ref{fig:grpo}, we introduce Group Relative Policy Optimization (GRPO). GRPO removes the need for additional value function estimation as required by PPO by instead adopting the mean reward across multiple responses to the same prompt as a baseline. Specifically, given a question $q$, GRPO samples a set of outputs $\{o_1, o_2, \cdots, o_G\}$ from $\pi_{\theta_{old}}$ and seeks to optimize the policy by maximizing:

\begin{equation}
\footnotesize
\begin{split}
    \mathcal{J}_{GRPO}(\theta) &= \mathbb{E}{[q \sim P(Q), \{o_i\}_{i=1}^G \sim \pi_{\theta_{old}}(O|q)]}  \\
    & \frac{1}{G}\sum_{i=1}^G\frac{1}{|o_i|} \sum_{t=1}^{|o_i|} \left\{ \min \left[ \frac{\pi_\theta(o_{i,t} | q, o_{i,<t})}{\pi_{\theta_{old}}(o_{i,t} | q, o_{i,<t})} \hat{A}_{i,t}, \text{clip} \left( \frac{\pi_\theta(o_{i,t} | q, o_{i,<t})}{\pi_{\theta_{old}}(o_{i,t} | q, o_{i,<t})}, 1 - \varepsilon, 1 + \varepsilon \right)  \hat{A}_{i,t} \right] - \beta \mathbb{D}_{KL}\left[\pi_{\theta} || \pi_{ref}\right]\right\} ,
\end{split}
\label{eq:GRPO-obj}
\end{equation}

with $\varepsilon$ and $\beta$ as hyper-parameters, and $\hat{A}_{i,t}$ representing an advantage that is determined based solely on the comparative rewards within each sampled output group (this process will be described in detail in the subsections below). By structuring advantage computation relatively across groups, GRPO harmonizes with the pairwise structure inherent in most reward model training datasets, which contain reward annotations based on comparisons of answers to identical queries. Furthermore, GRPO, instead of applying the KL regularization to the reward directly, incorporates it at the objective level by penalizing the KL divergence between the optimized and reference policies, simplifying the estimation of $\hat{A}_{i,t}$. Unlike the penalty formulation in (\ref{eq:PPO-reward}), we employ the following unbiased estimator for KL divergence as proposed in \citep{kl_approx}:

\begin{equation}
\small
    \mathbb{D}_{KL}\left[\pi_{\

\begin{algorithm}[t]
  \small
  \caption{Iterative Group Relative Policy Optimization}
  \textbf{Input} initial policy model $\pi_{\theta_{\text{init}}}$; reward models $r_\varphi$; task prompts $\mathcal{D}$; 
  hyperparameters $\varepsilon$, $\beta$, $\mu$
  \begin{algorithmic}[1]
    \State Initialize policy model $\pi_\theta \leftarrow \pi_{\theta_{\text{init}}}$
    \For{iteration = 1, \dots, I}
       \State Set reference model $\pi_{ref} \leftarrow \pi_{\theta}$
      \For{step = 1, \dots, M}
      \State Draw a batch $\mathcal{D}_b$ from $\mathcal{D}$
      \State Assign $\pi_{\theta_{old}} \leftarrow \pi_{\theta}$ 
      \State For each question $q \in \mathcal{D}_b$, generate $G$ responses $\{o_i\}_{i=1}^G \sim \pi_{\theta_{old}} (\cdot \mid q)$
      \State For every sampled output $o_i$, evaluate rewards $\{r_i\}_{i=1}^{G}$ using $r_{\varphi}$
      \State Use group-relative advantage estimation to calculate $\hat{A}_{i,t}$ for the $t$-th token in $o_i$
      \For{GRPO iteration = 1, \dots, $\mu$}
        \State Update $\pi_{\theta}$ by maximizing the GRPO objective as described in Equation \ref{eq:GC-GRPO}
      \EndFor
    \EndFor \State Enhance $r_\varphi$ via continual training employing a replay strategy.
    \EndFor 
  \end{algorithmic}
  \textbf{Output} $\pi_\theta$
  \label{alg:iter-grpo}
\end{algorithm}

\subsubsection{Outcome Supervision RL with GRPO} In formal terms, for each input question $q$, a set of outputs $\{o_1, o_2, \cdots, o_G\}$ is produced by sampling from the previous policy $\pi_{\theta_{old}}$. A reward model is applied to assign scores to each of these outputs, resulting in $G$ associated reward values $\mathbf{r} = \{r_1, r_2, \cdots, r_G\}$. These reward values are standardized by subtracting the mean of the group and dividing by their standard deviation. Within the framework of outcome supervision, this normalized reward is assigned at the final token of each output $o_i$ and used uniformly for all tokens in the sequence, setting the advantage as $\hat{A}_{i, t} = \widetilde{r}_i = \frac{r_i- {\rm mean}(\mathbf{r})}{{\rm std}(\mathbf{r})}$. The policy is then updated by optimizing the objective specified in equation (\ref{eq:GRPO-obj}).

\subsubsection{Process Supervision RL with GRPO} While outcome supervision provides a terminal reward per output, it may be inadequate for effectively guiding the learning process in challenging mathematical reasoning scenarios. Building on the approach in \cite{wang2023math}, process supervision is considered, which delivers feedback at each stage of the solution. Specifically, for question $q$ with $G$ generated outputs $\{o_1, o_2, \cdots, o_G\}$, a process reward model evaluates every intermediate step, assigning rewards as follows: $\mathbf{R} = \{ \{r_1^{index(1)},\cdots,r_1^{index(K_1)}\}, \cdots,  \{r_G^{index(1)},\cdots,r_G^{index(K_G)}\} \}$, where $index(j)$ denotes the terminal token of the $j$-th reasoning step, and $K_i$ is the number of steps in the $i$-th output. The stepwise rewards are normalized across the entire group by $\widetilde{r}_i^{index(j)} = \frac{r_i^{index(j)} - {\rm mean(\mathbf{R})}}{{\rm std(\mathbf{R})}}$.
Following this, process supervision determines the advantage for each token by summing the normalized rewards from that token’s step onwards: $\hat{A}_{i, t} = \sum_{index(j) \ge t} \widetilde{r}_i^{index(j)}$. The policy is subsequently optimized by maximizing the objective in equation (\ref{eq

\subsubsection{Iterative RL with GRPO} During reinforcement learning, the previously trained reward model may become inadequate for effectively guiding the latest iterations of the policy model. To address this, we implement iterative RL using GRPO. As detailed in Algorithm \ref{alg:iter-grpo}, the iterative GRPO process involves continually assembling new training datasets for the reward model from the outputs generated by the current policy model. The reward model is further trained with a replay buffer that maintains 10\% of its earlier data. Afterward, the updated policy model serves as the new reference model, and training of the policy model proceeds using the newly trained reward model.

\subsection{Training and Evaluating DeepSeekMath-RL}

RL training is carried out using DeepSeekMath-Instruct 7B. The RL dataset comprises chain-of-thought formatted items derived from GSM8K and MATH within the SFT data, totaling approximately 144K problems. We exclude unrelated SFT samples to examine how RL affects benchmark performance when no additional data is available during RL training. The method for assembling the reward model’s training dataset follows the approach described by \citep{wang2023math}. The initial reward model is trained on DeepSeekMath-Base 7B using a learning rate of 2e-5. For GRPO, the policy model employs a learning rate of 1e-6, and the KL penalty is set at 0.04. Each question is used to generate $64$ sampled completions. We set a maximum sequence length of 1024 and a batch size of 1024. After every exploration phase, the policy model undergoes a single update step. The evaluation of DeepSeekMath-RL 7B follows the protocols used for DeepSeekMath-Instruct 7B. Here, GSM8K and MATH with chain-of-thought reasoning are considered in-domain evaluation sets, whereas all remaining benchmarks are treated as out-of-domain tasks.

Table \ref{tab:sft_rl_math} presents a comparative analysis of both open- and closed-source models, evaluating their capabilities with chain-of-thought and tool-augmented reasoning on English and Chinese test suites. The results indicate that: 1) DeepSeekMath-RL 7B achieves 88.2\% accuracy on GSM8K and 51.7\% on MATH using chain-of-thought approaches, outperforming all open-source models within the 7B–70B parameter range and surpassing most closed-source counterparts. 2) Importantly, DeepSeekMath-RL 7B's training is limited to chain-of-thought instruction tuning data specifically from GSM8K and MATH, built upon DeepSeekMath-Instruct 7B. Despite its narrower training corpus, it demonstrates superior outcomes compared to DeepSeekMath-Instruct 7B across all benchmarks, highlighting the benefits gained from reinforcement learning.

\section{Discussion}
In this section, we present insights obtained from both our pre-training and RL experiments.

\subsection{Insights from Pre-Training}

We begin by detailing our experiences during the pre-training phase. Unless noted otherwise, the training configurations are as specified in Section \ref{sec:quality-policy}. Additionally, the term DeepSeekMath Corpus in this context refers to an 89B-token dataset generated during the second cycle of our data gathering process.

\subsubsection{Code Training Enhances Mathematical Reasoning}

It is commonly hypothesized—though not empirically validated—that code-related pre-training bolsters models' reasoning capabilities. Here, we attempt to provide partial evidence supporting this, particularly in mathematical tasks: training on code data enhances mathematical reasoning skills, with or without reliance on tools.

To investigate the influence of code training on mathematical reasoning, we conducted experiments using both a two-stage and a single-stage training protocol:

\textbf{Two-Stage Training}

\begin{itemize}[topsep=0pt]
    \item \textbf{Code Training for 400B Tokens $\rightarrow$ Math Training for 150B Tokens}: We subject DeepSeek-LLM 1.3B to an initial phase involving 400B code tokens, then proceed with an additional 150B math tokens for continued training;
    \item \textbf{General Training for 400B Tokens $\rightarrow$ Math Training for 150B Tokens}: As a baseline for comparison, a parallel experiment utilizes a dataset of 400B general tokens (sourced from DeepSeek-AI's comprehensive general corpus) for the first training stage, rather than code tokens. This is designed to assess whether code-centric pretraining confers any distinct advantage in enhancing the model's mathematical reasoning abilities over generic language pretraining.
\end{itemize}

\textbf{One-Stage Training}

\begin{itemize}[topsep=0pt]
    \item \textbf{Math Training for 150B Tokens}: We train DeepSeek-LLM 1.3B exclusively with 150B tokens comprising mathematical content;
    \item \textbf{Training on a mixture of 400B Code Tokens and 150B Math Tokens}: Introducing math-focused training directly after code training has a detrimental effect on coding proficiency. To probe this phenomenon, we conduct an alternative training regimen in which code tokens and math tokens are combined and presented jointly in a single-stage setup. The goal is to evaluate whether this mixture can simultaneously sustain improvements in mathematical reasoning and mitigate the adverse effects of catastrophic forgetting.
\end{itemize}

\paragraph{Results}
The experimental outcomes under these various training paradigms are presented in Table \ref{tab:code-math} and Table \ref{tab:code-math-general-eval}.

Across both the two-stage and one-stage protocols, prior exposure to code data consistently enhances program-assisted mathematical problem-solving. As indicated by Table \ref{tab:code-math}, in the two-stage regime, training with code data alone leads to marked improvements in the model's capability to resolve problems in datasets such as GSM8K and MATH using Python code. Supplementary math-specific training during the second stage contributes additional gains. Notably, when adopting the one-stage approach, jointly presenting code and math tokens significantly curtails the issue of catastrophic forgetting commonly encountered in sequential two-stage setups. This joint strategy also enhances performance on both coding (Table \ref{tab:code-math-general-eval}) and program-augmented mathematical tasks (Table \ref{tab:code-math}).

Furthermore, training on code tokens benefits mathematical reasoning tasks that do not involve external tool use. Within the two-stage context, code training as the initial phase provides a moderate performance boost on these tasks. It also seems to facilitate the subsequent math training, producing the highest results overall. Conversely, blending code and math tokens in a unified one-stage process negatively impacts mathematical reasoning in tool-free scenarios. One plausible explanation is that the DeepSeek-LLM 1.3B model may not possess sufficient parameter capacity to internalize both code-based and mathematical information in tandem when presented together.

\subsubsection{Limited Effect of ArXiv Papers on Mathematical Reasoning}

Although the inclusion of arXiv papers as part of mathematical pre-training corpora has become standard practice \citep{minerva,gpt-f,llemma,mathpile}, systematic evaluation of their contribution to mathematical reasoning ability remains sparse. Somewhat unexpectedly, our results suggest that pre-training with arXiv content yields little to no enhancement in mathematical reasoning skills. We evaluated this hypothesis across multiple model configurations—specifically DeepSeek-LLM 1.3B and DeepSeek-Coder-Base-v1.5 7B \citep{deepseek-coder}—using several forms of arXiv-derived datasets processed by distinct pipelines:

\begin{itemize}[topsep=0pt]
    \item \textbf{MathPile} \citep{mathpile}: A collection consisting of 8.9 billion tokens, generated through data cleaning and heuristic-based filtering, with scientific arXiv papers comprising over 85\% of its content;
    \item \textbf{ArXiv-RedPajama} \citep{redpajama}: This dataset includes all available arXiv LaTeX source files after eliminating preambles, comments, macro definitions, and bibliographies, containing 28.0 billion tokens in total.
\end{itemize}

For these experiments, we separately trained DeepSeek-LLM 1.3B on 150B tokens and DeepSeek-Coder-Base-v1.5 7B on 40B tokens for each of the above arXiv corpora. The findings consistently indicate that models pre-trained solely with arXiv data do not demonstrate significant improvements and occasionally even regress in performance across a wide range of mathematical benchmarks. The evaluated benchmarks—encompassing datasets for quantitative reasoning such as GSM8K and MATH (see Table \ref{tab:arxiv-cot}), standardized multiple-choice problems such as those in MMLU-STEM (Table \ref{tab:arxiv-cot}), as well as formal mathematics like miniF2F (Table \ref{tab:arxiv_atp})—showed this trend irrespective of problem difficulty.

Nevertheless, several caveats temper these findings, and they should not be considered definitive. Key limitations in the current investigation include:

\begin{itemize}[topsep=0pt]
    \item The unexplored effect of arXiv data on niche mathematical applications not assessed in this study, such as the conversion of formal proofs or theorems to informal narrative (the informalization task);
    \item The potential interaction effects of arXiv material when jointly used with other types of pre-training data;
    \item The possibility that advantages from arXiv data could emerge with significantly larger models.
\end{itemize}

Hence, there is ample opportunity for future work to further clarify these questions.

\subsection{Reinforcement Learning: Key Insights}
\subsubsection{A Unified Framework} We now introduce a comprehensive perspective for analyzing several training algorithms, including SFT, RFT, DPO, PPO, and GRPO, and we design experiments to probe the critical aspects within this paradigm. The general form of the parameter gradient $\theta$ for any such training strategy can be expressed as:

\begin{equation}
    \nabla_{\theta}\mathcal{J}_{\textcolor{red}{\mathcal{A}}}(\theta) = \mathbb{E}[\underbrace{(q,o) \sim \textcolor{red}{\mathcal{D}}}_{Data \ Source}]\left( \frac{1}{|o|} \sum_{t=1}^{|o|}  \underbrace{GC_{{\mathcal{A}}}(q, o, t, \textcolor{red}{\pi_{{rf}}})}_{Gradient \ Coefficient}  \nabla_{\theta}\log \pi_{\theta}(o_t | q, o_{<t})\right).
\label{eq:objective}
\end{equation}

Three principal elements are present in this formulation: 1) \textit{Data Source $\mathcal{D}$}, which specifies the training dataset utilized; 2) \textit{Reward Function $\pi_{{rf}}$}, which serves as the provider of reward signals guiding the learning process; and 3) \textit{Algorithm $\mathcal{A}$}, which integrates both the dataset and reward signals to compute the gradient coefficient $GC$, thereby influencing the degree of updating—either penalization or reinforcement—applied to each example. We provide an analysis of several commonly-used methods through this unified perspective:

\begin{itemize}
    \item  \textbf{Supervised Fine-tuning (SFT)}: SFT adapts the pretrained model through additional training on datasets curated by humans.
    \item \textbf{Rejection Sampling Fine-tuning (RFT)}: RFT enhances the SFT model by further fine-tuning it with outputs—generated by the SFT model in response to the same questions—which have been filtered based on answer accuracy.
    \item \textbf{Direct Preference Optimization (DPO)}: DPO applies further refinement by training the SFT model on expanded datasets containing sampled outputs, leveraging a pairwise preference-based loss.
    \item \textbf{Online Rejection Sampling Fine-tuning (Online RFT)}: In contrast with RFT, Online RFT employs the SFT model for initial policy and incrementally fine-tunes the model using newly generated outputs from the evolving policy model during training.
    \item \textbf{PPO/GRPO}: PPO/GRPO begins with the SFT model as the base policy and then further updates the model through reinforcement learning using outputs produced by the current policy in real-time.
\end{itemize}

We present a summary of the constituent elements of these approaches in Table \ref{tab:data_policy}. For an in-depth exposition of the derivation process, see Appendix \ref{app:analysis-rl}.

\paragraph{Reflection on Data Source} The categorization of data sources is bifurcated into online sampling and offline sampling. Online sampling indicates that the dataset utilized during training arises from the contemporaneous exploration performed by the current policy model. Conversely, offline sampling involves using data gathered from the initial SFT model’s sampling process. RFT and DPO exemplify methods reliant on offline data, whereas Online RFT and GRPO operate with data acquired online.

According to Figure \ref{fig:rl-analysis}, our empirical results demonstrate that Online RFT markedly surpasses RFT across both benchmarks considered. More precisely, Online RFT’s performance matches that of RFT in early training but secures a distinct and growing lead as training proceeds, underscoring the effectiveness of the online training paradigm. This observation aligns with expectations: initially, the actor closely resembles the SFT model, yielding sampled data with minimal discrepancies. However, as training advances, actor-generated samples deviate more substantially, thereby amplifying the benefits of utilizing freshly sampled data in real time.

\paragraph{Reflection on Gradient Coefficient} The transformation of the reward signal into a gradient coefficient, which subsequently guides parameter updates, constitutes a core step in the algorithm. In our experimentation, the reward function is split into two forms: `Rule' and `Model.' The `Rule' approach assesses responses based on their correctness, while the `Model' approach entails developing a dedicated reward model that assigns a score to each output; this reward model is itself trained on data generated using the rule-based approach. Equations \ref{eq:GC-RFT} and \ref{eq:GC-GRPO} draw attention to a critical distinction between GRPO and Online RFT: specifically, GRPO modulates its gradient coefficient in accordance with the reward model’s output, which enables variable positive or negative updates depending on response quality. In contrast, Online RFT does not make such distinctions—incorrect responses are not penalized, and correct responses all receive the same degree of reinforcement.

As illustrated in Figure \ref{fig:rl-analysis}, GRPO outperforms online RFT, underscoring the effectiveness of adjusting the coefficients for positive and negative gradients. Moreover, GRPO+PS demonstrates higher effectiveness relative to GRPO+OS, which underscores the advantage of employing fine-grained, step-dependent gradient adjustments. We also investigate iterative RL procedures, implementing two cycles in our experiments. According to Figure \ref{fig:iter-rl}, it is evident that iterative RL considerably boosts performance, particularly during the initial iteration.

\subsubsection{Why RL Works?} In this study, reinforcement learning is carried out on a selected portion of the instruction tuning dataset, and this approach yields marked performance gains over the instruction-tuned baseline. To clarify the reasons behind the efficacy of reinforcement learning, we analyze both Pass@K and Maj@K accuracy metrics for the Instruct and RL models across two benchmarks. As depicted in Figure \ref{fig:maj-pass}, reinforcement learning raises Maj@K accuracy, but does not affect Pass@K. These observations suggest that reinforcement learning increases overall robustness in the output distribution, i.e., \textbf{the observed performance gains arise from an increased likelihood of the correct answer appearing in the TopK outputs, rather than from fundamental improvements in the model’s abilities.} Analogous findings were reported in \citep{wang2023making}, which identified a \textbf{misalignment problem} affecting reasoning tasks in SFT models, and showed that aligning model preferences through targeted strategies can improve reasoning performance \citep{yuan2023rrhf,song2023preference,wang2023making}.

\subsubsection{How to Achieve More Effective RL?}
\label{sec:effec_rl}
Our results verify that reinforcement learning proves highly effective for mathematical reasoning applications. We additionally present a cohesive framework for interpreting prominent training paradigms, categorizing all such methods as variants or simplifications of reinforcement learning. Equation \ref{eq:objective} encapsulates three central elements of this framework: Data Source, Algorithm, and Reward Function. We outline several prospective research avenues within each of these components.

\paragraph{Data Source} The data source underpins every training approach. For RL-based approaches, this refers specifically to unlabeled problems paired with model-generated outputs. In our study, only question data from the instruction tuning phase was used, and response samples were generated via simple nucleus sampling. We suspect this design choice partially accounts for why our RL approach selectively enhanced Maj@K scores. Moving forward, we plan to extend our RL pipeline to tackle prompts from out-of-distribution sources, while adopting \textbf{advanced sampling (decoding) methods}, such as tree-search algorithms \citep{yao2023tree}. In addition, the adoption of \textbf{efficient inference strategies} \citep{xia-etal-2023-speculative,leviathan2023fast,kwon2023efficient,xia2024unlocking}, which crucially shape the exploratory behavior of policy models, will be integral to further advancements.

\paragraph{Algorithms} Algorithms utilize both data and reward signals to compute gradients, thereby guiding updates to model parameters. According to Equation \ref{eq:objective}, most current approaches essentially \textbf{TRUST} the reward function’s feedback when adjusting the conditional probability associated with particular tokens. Nonetheless, there is no guarantee that the reward signal will remain accurate in all circumstances, particularly within highly complex domains. As an illustrative case, even datasets such as PRM800K \citep{lightman2023let}, which have been meticulously labeled by experienced annotators, still present an error rate close to 20\% in their annotations\footnote{\url{https://github.com/openai/prm800k/issues/12\#issuecomment-1728491852}}. Consequently, our focus will be on reinforcement learning algorithms designed to be resilient to unreliable or noisy rewards. We anticipate that alignment techniques following the \textbf{WEAK-TO-STRONG} framework \citep{burns2023weak} could fundamentally reshape existing approaches to model training.

\paragraph{Reward Function} The reward function serves as the principal origin of the learning signal in reinforcement learning and is often instantiated via a neural reward model. We identify three key research avenues for these models: 1) \textbf{Improving the generalization capacity of the reward model.} It is crucial for reward models to generalize beyond the training distribution and cope with diverse and sophisticated outputs; failing this, reinforcement learning may merely perpetuate the existing behavior of LLMs, rather than enhancing their intrinsic abilities; 2) \textbf{Quantifying the reward model's uncertainty.} Accounting for this uncertainty might be instrumental in bridging the gap between less robust reward models and the weak-to-strong alignment methods; 3) \textbf{Constructing process reward models of high quality and efficiency} that are capable of supplying detailed, process-level training feedback, particularly in domains requiring stepwise reasoning \citep{lightman2023let,wang2023math}.

\section{Conclusion, Limitation, and Future Work} In this work, we introduce DeepSeekMath, which surpasses all other open-source models on the competitive MATH benchmark and attains near-parity with proprietary systems. Built upon the DeepSeek-Coder-v1.5 7B foundation, DeepSeekMath undergoes sustained training for 500B tokens, drawing extensively on a corpus that includes 120B math tokens primarily collected from Common Crawl. Through a comprehensive ablation study, we determine that web-derived resources provide substantial high-quality mathematical content, whereas arXiv does not contribute as significantly as anticipated. We present Group Relative Policy Optimization (GRPO), a modified Proximal Policy Optimization (PPO) approach, which enhances mathematical reasoning abilities while requiring less memory. Our experimental analysis demonstrates the efficacy of GRPO even when DeepSeekMath-Instruct 7B attains strong performance on standard benchmarks. Additionally, we propose a unifying perspective for interpreting several related techniques and highlight future directions to achieve more effective reinforcement learning.

Despite achieving notable results on quantitative reasoning tasks, DeepSeekMath's proficiency in geometry and formal theorem-proving remains behind that of leading closed models. Specifically, during internal testing, the model struggled with triangle and ellipse-related questions, suggesting potential bias in both pre-training and fine-tuning dataset selection. Furthermore, constrained by model size, DeepSeekMath trails GPT-4 in few-shot scenarios; whereas GPT-4's accuracy improves with few-shot prompting, DeepSeekMath's zero-shot and few-shot results remain comparable. Moving forward, our objective is to refine the engineered data selection pipeline for constructing a higher-quality pre-training corpus. We also intend to pursue the avenues discussed in Section \ref{sec:effec_rl} for further advancements in reinforcement learning techniques for LLMs.

\end{document}